\textbf{Parametric Resonance}
The system satisfies the equation of motion:
\[
\ddot{x} + \omega (t) x = 0
\]
where $\omega (t)$ is a periodic function with period $T$.

\textbf{Conditions for Parametric Resonance}
We substitute the solution space of the motion and obtain a set of solutions. Due to the periodicity of the equation, we obtain a solution in the form:
\[
x_i (t + nT) = \mu_i^n x_i
\]
where the solution can be expressed as a product of periodic functions. This leads to the condition:
\[
\mu_1 \mu_2 = 1 \Rightarrow x_2 x_1 - x_1 x_2 = \text{const}
\]
Thus, parametric resonance occurs if $\mu_1 \neq 1$.

\textbf{Growth of Perturbation for Mathieu's Equation}
In general, when does the perturbation $\omega (t)$ grow?  
The approximate solution is given by:
\[
\omega_0^2 (1 + \varepsilon \cos \gamma t)
\]
For small perturbations, we approximate $\gamma$ as:
\[
\gamma = \omega_0 + \delta
\]
where the solution takes the form:
\[
a(t) \cos \left(\omega_0 + \frac{\delta}{2} \right)t + b(t) \sin \left(\omega_0 + \frac{\delta}{2} \right)t
\]
where $a, b$ are evaluated at points of minimal perturbation. Under certain conditions, we obtain the constraint:
\[
|\delta| < \frac{1}{2} \omega_0^2 \varepsilon \frac{\epsilon^2}{32} \omega_0^2.
\]

\subsection*{Trick for Finding $\eta$}
We start with the perturbative equation in the form:
\[
\dot{\phi} = -\frac{\partial V}{\partial \phi} + f (\phi, t)
\]
where $f_1$ represents the weak force on the motion. We solve for the slow-time solution:
\[
\phi = \Phi + \eta
\]
where the effective potential is given by:
\[
V_{\text{eff}} (\Phi) = V (\Phi) + \frac{1}{2} \langle \eta^2 \rangle.
\]

\subsection*{Scattering - Differential Cross Section}
The inclusion of the solid angle for radial coordinates is given by:
\[
d\Omega = \frac{dS}{r^2} = \sin \theta d\theta d\varphi
\]
which integrates as:
\[
d\Omega = 2\pi \sin \theta d\theta.
\]

\subsection*{Scattering Cross Section}
The scattered area distribution in terms of the solid angle is given by:
\[
\frac{\partial \sigma}{\partial \Omega} = \frac{\rho (\theta)}{\sin \theta} \left| \frac{d\rho}{d\theta} \right|
\]

\subsection*{Total Cross Section}
\[
\sigma_T = \int \frac{d\sigma}{d\Omega} d\Omega = 2\pi \int_0^\pi \rho (\theta) \left| \frac{d\rho}{d\theta} \right| d\theta.
\]

\subsection*{Central Potential}
We find that the total energy is:
\[
E = \frac{1}{2} \mu \omega_0^2 \rho^2
\]
which leads to the relation:
\[
\rho^2 \left(1 - \frac{V (r_{\min})}{E} \right) = \rho^2.
\]
Thus, we establish the relation between $\theta$ in the central potential.
